\begin{apendicesenv}

% Imprime uma página indicando o início dos apêndices
\partapendices

% ----------------------------------------------------------
\chapter{Prompts utilizados na Engenharia de Prompt}
\label{ap:prompts}
% ----------------------------------------------------------

Este apêndice reproduz, na íntegra e exatamente como foram enviados aos modelos, os \textit{prompts} de cada nível de normalização (N1, N2 e N3). Cada \textit{prompt} reside em um arquivo de texto próprio no diretório \texttt{prompts/} do repositório. Os textos são apresentados sem acentuação nas instruções, tal como nos arquivos-fonte; a acentuação aparece apenas nos exemplos de saída esperada. Os marcadores entre chaves --- \verb|{text}|, \verb|{text_n1}|, \verb|{aplicabilidade}|, \verb|{requisito}|, \verb|{selecao}| e \verb|{execao}| --- são substituídos, em tempo de execução, pelo texto bruto da norma, pela regra atômica do N1 e pelo operador N2 correspondente.

% ----------------------------------------------------------
\section{N1 --- \textit{Prompt Direct}}
% ----------------------------------------------------------

\begin{lstlisting}[caption={\textit{Prompt} do nível N1 (\textit{Prompt Direct}) --- \texttt{prompts/n1.txt}.}, label={lst:n1}]
Voce e um reescritor. Transforme o texto em sentencas RASE N1.

Regras:
1) Quebre em sentencas curtas e diretas.
2) Cada sentenca deve ter uma unica regra computavel.
3) Nao invente elementos (aplicabilidade, selecao, requisito, excecao). Preserve apenas o que existir.
4) Nao adicione explicacoes, titulos, bullets, numeracao ou o texto original.
5) Saida: apenas as sentencas, uma por linha, terminadas com ponto final.

Exemplo:
Entrada:
A inclinacao transversal da superficie deve ser de ate 2 % para pisos internos e de ate 3 % para pisos externos.
Saida:
Pisos internos devem ter inclinacao transversal de no maximo 2%.
Pisos externos devem ter inclinacao transversal de no maximo 3%.

TEXTO_INICIO
{text}
TEXTO_FIM
\end{lstlisting}

% ----------------------------------------------------------
\section{N2 --- \textit{Chain-of-Thought}}
% ----------------------------------------------------------

\begin{lstlisting}[caption={\textit{Prompt} do nível N2 (\textit{Chain-of-Thought}) --- \texttt{prompts/n2.txt}.}, label={lst:n2}]
Extrator RASE N2.
Use APENAS o Texto N1 para extrair os elementos. O Texto completo e apenas referencia.

Regras (ordem fixa):
1) aplicabilidade (opcional): onde/quando se aplica, sem verbos.
2) selecao (opcional): subconjunto da aplicabilidade, sem verbos.
3) execao (opcional): casos que NAO seguem a regra.
4) requisito (obrigatorio): acao/condicao principal, comeca com verbo.

Regras de saida:
- Retorne exatamente 4 linhas no formato abaixo.
- Cada campo deve aparecer no maximo uma vez.
- Se nao existir, use "" (string vazia).
- Nao adicione explicacoes, listas ou texto extra.

Exemplo (formato):
Texto completo:
"As areas ... Norma."
Texto N1:
"As areas de qualquer espaco ou edificacao de uso publico ou coletivo devem ser servidas de uma ou mais rotas acessiveis."
Resposta (4 linhas):
aplicabilidade: As areas de qualquer espaco ou edificacao
selecao: uso publico ou coletivo
execao: ""
requisito: devem ser servidas de uma ou mais rotas acessiveis

Agora processe:
Texto completo:
"{text}"
Texto N1:
"{text_n1}"

Resposta (4 linhas):
aplicabilidade:
selecao:
execao:
requisito:
\end{lstlisting}

% ----------------------------------------------------------
\section{N3 --- \textit{Few-Shot}}
% ----------------------------------------------------------

O nível N3 emprega um \textit{prompt} por operador RASE. Os quatro compartilham a mesma estrutura --- esquema JSON de saída, heurísticas de mapeamento e exemplos \textit{few-shot} ---, variando apenas as heurísticas e os exemplos próprios de cada operador.

% ----------------------------------------------------------
\subsection{Aplicabilidade}
% ----------------------------------------------------------

\begin{lstlisting}[caption={\textit{Prompt} do operador Aplicabilidade (N3) --- \texttt{prompts/n3\_aplicabilidade.txt}.}, label={lst:n3-apl}]
Extrator RASE N3 (aplicabilidade).
Use APENAS o operador N2 de aplicabilidade para gerar as propriedades N3. O texto completo/N1 e apenas referencia, se existir.

Regras:
1) Gere um JSON com os campos: type, object, property, comparation, target, unit.
2) type deve ser "aplicabilidade". Se o operador estiver vazio, todos os campos devem ser "".
3) Use o operador como fonte principal; use Texto completo/N1 apenas para desambiguar.
4) object: substantivo(s) principal(is) do trecho; se houver mais de um, separe por "; ".
5) property/comparation/target:
   - Quando o trecho indicar uso (ex: "edificacoes residenciais"), use property "uso".
     - comparation "=" quando o texto definir o uso.
     - comparation "contém" quando for subconjunto/parte do uso.
   - Quando o trecho indicar "acessos", use object "componente", property "tipo",
     comparation "inclui", target "porta; passagem".
   - Quando o trecho indicar "edificacoes e equipamentos urbanos", use
     object "edificação; equipamento urbano" e target "obra nova".
   - Quando o trecho indicar tipo do componente (ex: entradas, rotas), use property "tipo", comparation "inclui".
   - Quando nao houver propriedade clara, deixe property/comparation/target vazios.
6) unit: use "" (vazio) a menos que exista unidade explicita.

Regras de saida:
- Retorne somente um JSON valido.
- Nao adicione explicacoes ou texto extra.

Exemplo:
Operador N2 (aplicabilidade):
As edificacoes residenciais multifamiliares, condominios e conjuntos habitacionais

Resposta (JSON):
{
  "type": "aplicabilidade",
  "object": "edificação",
  "property": "uso",
  "comparation": "=",
  "target": "residenciais multifamiliares, condomínios e conjuntos habitacionais",
  "unit": ""
}

Exemplo:
Operador N2 (aplicabilidade):
Os acessos

Resposta (JSON):
{
  "type": "aplicabilidade",
  "object": "componente",
  "property": "tipo",
  "comparation": "inclui",
  "target": "porta; passagem",
  "unit": ""
}

Exemplo:
Operador N2 (aplicabilidade):
edificacoes e equipamentos urbanos

Resposta (JSON):
{
  "type": "aplicabilidade",
  "object": "edificação; equipamento urbano",
  "property": "",
  "comparation": "",
  "target": "obra nova",
  "unit": ""
}

Exemplo:
Operador N2 (aplicabilidade):
As unidades autonomas acessiveis

Resposta (JSON):
{
  "type": "aplicabilidade",
  "object": "espaço",
  "property": "uso",
  "comparation": "contém",
  "target": "comum",
  "unit": ""
}

Agora processe:
Texto completo:
"{text}"
Texto N1:
"{text_n1}"
Operador N2 (aplicabilidade):
"{aplicabilidade}"

Resposta (JSON):
\end{lstlisting}

% ----------------------------------------------------------
\subsection{Requisito}
% ----------------------------------------------------------

\begin{lstlisting}[caption={\textit{Prompt} do operador Requisito (N3) --- \texttt{prompts/n3\_requisito.txt}.}, label={lst:n3-req}]
Extrator RASE N3 (requisito).
Use APENAS o operador N2 de requisito para gerar as propriedades N3. O texto completo/N1 e apenas referencia, se existir.

Regras:
1) Gere um JSON com os campos: type, object, property, comparation, target, unit.
2) type deve ser "requisito". Se o operador estiver vazio, todos os campos devem ser "".
3) Use o operador como fonte principal; use Texto completo/N1 apenas para desambiguar.
4) object: substantivo(s) principal(is) do trecho; se houver mais de um, separe por "; ".
5) property/comparation/target:
   - "acessivel" -> property "acessível", comparation "=", target "VERDADEIRO"
   - "conectadas as rotas acessiveis" / "vinculados atraves de rota acessivel" ->
     property "conectado a rota acessível" (ou "conectado à rota acessível" se houver crase),
     comparation "=", target "VERDADEIRO"
   - "livres de obstaculos" -> property "possui obstáculo", comparation "=", target "FALSO"
   - "devem ser servidas de rota(s) acessiveis" -> property "uso", comparation "=", target "VERDADEIRO"
   - Quando nao houver propriedade clara, deixe property/comparation/target vazios.
6) unit: use "" (vazio) a menos que exista unidade explicita.

Regras de saida:
- Retorne somente um JSON valido.
- Nao adicione explicacoes ou texto extra.

Exemplo:
Operador N2 (requisito):
devem permanecer livres de quaisquer obstaculos de forma permanente

Resposta (JSON):
{
  "type": "requisito",
  "object": "porta; passagem",
  "property": "possui obstáculo",
  "comparation": "=",
  "target": "FALSO",
  "unit": ""
}

Exemplo:
Operador N2 (requisito):
devem ser acessiveis

Resposta (JSON):
{
  "type": "requisito",
  "object": "porta; espaço",
  "property": "acessível",
  "comparation": "=",
  "target": "VERDADEIRO",
  "unit": ""
}

Exemplo:
Operador N2 (requisito):
devem ser servidas de uma ou mais rotas acessiveis

Resposta (JSON):
{
  "type": "requisito",
  "object": "espaço ou edificação",
  "property": "uso",
  "comparation": "=",
  "target": "VERDADEIRO",
  "unit": ""
}

Agora processe:
Texto completo:
"{text}"
Texto N1:
"{text_n1}"
Operador N2 (requisito):
"{requisito}"

Resposta (JSON):
\end{lstlisting}

% ----------------------------------------------------------
\subsection{Seleção}
% ----------------------------------------------------------

\begin{lstlisting}[caption={\textit{Prompt} do operador Seleção (N3) --- \texttt{prompts/n3\_selecao.txt}.}, label={lst:n3-sel}]
Extrator RASE N3 (selecao).
Use APENAS o operador N2 de selecao para gerar as propriedades N3. O texto completo/N1 e apenas referencia, se existir.

Regras:
1) Gere um JSON com os campos: type, object, property, comparation, target, unit.
2) type deve ser "selecao". Se o operador estiver vazio, todos os campos devem ser "".
3) Use o operador como fonte principal; use Texto completo/N1 apenas para desambiguar.
4) object: substantivo(s) principal(is) do trecho; se houver mais de um, separe por "; ".
5) property/comparation/target:
   - "uso publico/coletivo/comum/restrito" -> property "uso".
     - comparation "=" quando o texto definir o uso.
     - comparation "contém" quando o texto indicar subconjunto (ex: "uso comum", "uso restrito").
   - listagem de tipos (entradas, rotas, circulacoes) -> property "tipo", comparation "inclui",
     target com tipos normalizados (ex: "porta; passagem; circulação").
   - Quando a selecao for apenas um qualificador (ex: "internos", "externos", "eventuais"),
     deixe todos os campos vazios.
   - Quando nao houver propriedade clara, deixe property/comparation/target vazios.
6) unit: use "" (vazio) a menos que exista unidade explicita.

Regras de saida:
- Retorne somente um JSON valido.
- Nao adicione explicacoes ou texto extra.

Exemplo:
Operador N2 (selecao):
uso publico ou coletivo

Resposta (JSON):
{
  "type": "selecao",
  "object": "espaço",
  "property": "uso",
  "comparation": "=",
  "target": "público ou coletivo",
  "unit": ""
}

Exemplo:
Operador N2 (selecao):
todas as entradas, bem como as rotas de interligacao as funcoes do edificio

Resposta (JSON):
{
  "type": "selecao",
  "object": "componente; espaço",
  "property": "tipo",
  "comparation": "inclui",
  "target": "porta; passagem; circulação",
  "unit": ""
}

Agora processe:
Texto completo:
"{text}"
Texto N1:
"{text_n1}"
Operador N2 (selecao):
"{selecao}"

Resposta (JSON):
\end{lstlisting}

% ----------------------------------------------------------
\subsection{Exceção}
% ----------------------------------------------------------

\begin{lstlisting}[caption={\textit{Prompt} do operador Exceção (N3) --- \texttt{prompts/n3\_execao.txt}.}, label={lst:n3-exc}]
Extrator RASE N3 (execao).
Use APENAS o operador N2 de execao para gerar as propriedades N3. O texto completo/N1 e apenas referencia, se existir.

Regras:
1) Gere um JSON com os campos: type, object, property, comparation, target, unit.
2) type deve ser "excecao". Se o operador estiver vazio, todos os campos devem ser "".
3) Use o operador como fonte principal; use Texto completo/N1 apenas para desambiguar.
4) object: substantivo(s) principal(is) do trecho; se houver mais de um, separe por "; ".
5) property/comparation/target:
   - "uso restrito" ou listagem de areas excluidas -> property "uso", comparation "contém",
     target com lista normalizada (ex: "casa de máquinas, barrilete e técnico").
   - Quando for apenas condicao/tempo ("em reformas", "se ..."), deixe todos os campos vazios.
   - Quando nao houver propriedade clara, deixe property/comparation/target vazios.
6) unit: use "" (vazio) a menos que exista unidade explicita.

Regras de saida:
- Retorne somente um JSON valido.
- Nao adicione explicacoes ou texto extra.

Exemplo:
Operador N2 (execao):
Areas de uso restrito, como casas de maquinas, barriletes e passagem de uso tecnico

Resposta (JSON):
{
  "type": "excecao",
  "object": "espaço",
  "property": "uso",
  "comparation": "contém",
  "target": "casa de máquinas, barrilete e técnico",
  "unit": ""
}

Exemplo:
Operador N2 (execao):
Em reformas

Resposta (JSON):
{
  "type": "",
  "object": "",
  "property": "",
  "comparation": "",
  "target": "",
  "unit": ""
}

Agora processe:
Texto completo:
"{text}"
Texto N1:
"{text_n1}"
Operador N2 (execao):
"{execao}"

Resposta (JSON):
\end{lstlisting}

\end{apendicesenv}
